\item Un parámetro común que puede utilizarse para predecir turbulencia en el flujo de un fluido es el llamado número de Reynolds. El número de Reynolds para el flujo de un fluido en una tubería es una cantidad adimensional definida por
$$ \text{Re} = \frac{\rho v d}{\mu} $$
donde $\rho$ es la densidad del fluido, $v$ es su rapidez, $d$ es el diámetro interno de la tubería y $\mu$ es la viscosidad del fluido. La viscosidad es una medida de la resistencia interna de un líquido al flujo y tiene unidades de $\text{Pa}\cdot\text{s}$. El criterio para el tipo de flujo es como sigue:
\\begin{itemize}
    \\item Si $\\text{Re} < 2300$, el flujo es laminar.
    \\item Si $2300 < \\text{Re} < 4000$, el flujo está en una región de transición entre laminar y turbulento.
    \\item Si $\\text{Re} > 4000$, el flujo es turbulento.
\\end{itemize}
(a) Sangre de densidad $1.06 \times 10^3 \text{ kg/m}^3$ y viscosidad $3.00 \times 10^{-3} \text{ Pa}\cdot\text{s}$ se modela como un líquido puro, es decir, se ignora la presencia de los glóbulos rojos. Suponga que la sangre fluye en una gran arteria de radio $1.50 \text{ cm}$ con una rapidez de $0.067\,0 \text{ m/s}$. Demuestre que el flujo es laminar. (b) Imagine que la arteria finaliza en un solo capilar, así que el radio de la arteria se reduce a un valor mucho más pequeño. ¿Cuál es el radio del capilar que provocaría un flujo turbulento? (c) Los capilares reales tienen radios de aproximadamente $5$-$10 \text{ micrómetros}$, mucho más pequeños que el valor en el inciso (b). ¿Por qué el flujo en los capilares reales no se hace turbulento?
